\documentclass[11pt]{article}
\usepackage[margin=0.7in]{geometry}
\usepackage{amsmath,amssymb}
\usepackage{enumitem}
\usepackage{tikz}
\usepackage{pgfplots}
\pgfplotsset{compat=1.18}
\usepackage{fancyhdr}
\usepackage{parskip}
\setlist[enumerate]{leftmargin=*, itemsep=6pt}
\pagestyle{fancy}
\fancyhf{}
\fancyfoot[L]{\small\textit{en}Vision\textsuperscript{\textregistered} Algebra 2}
\fancyfoot[C]{\small 3-3 Lesson Quiz, short form}
\fancyfoot[R]{\small\thepage}
\renewcommand{\headrulewidth}{0pt}
\newcommand{\choice}[1]{\tikz[baseline=-0.6ex]\node[draw, circle, inner sep=0.6pt, minimum size=3.3mm, font=\scriptsize\bfseries] {#1};\;}
\newcommand{\checkbox}{$\square$\;}

\begin{document}

\noindent\textbf{Name}\; \rule{6cm}{0.4pt} \hfill Period \rule{1.2cm}{0.4pt} \quad Date \rule{2.2cm}{0.4pt}

\vspace{4pt}\noindent\rule{\textwidth}{0.6pt}

\vspace{2pt}\noindent{\LARGE\textbf{3-3 \; Lesson Quiz: Polynomial Identities}}

\noindent{\small short form \hfill questions keep their Savvas numbers (2, 4, 5); show your work}

\vspace{2pt}\noindent\rule{\textwidth}{0.6pt}

\vspace{6pt}
\textbf{2.} \textbf{Part A} Use polynomial identities to multiply $(5-4x^3)(5+4x^3)$.

\begin{minipage}[t]{0.48\textwidth}
\begin{itemize}[leftmargin=1.2em, itemsep=5pt]
\item[] \choice{A} $25-4x^9$
\item[] \choice{B} $25-40x^3+16x^6$
\item[] \choice{C} $25-4x^6$
\item[] \choice{D} $25-16x^6$
\end{itemize}
\end{minipage}\hfill
\begin{minipage}[t]{0.48\textwidth}
\textbf{Part B} What identity did you use?
\begin{itemize}[leftmargin=1.2em, itemsep=5pt]
\item[] \choice{A} Difference of squares
\item[] \choice{B} Difference of cubes
\item[] \choice{C} Sum of cubes
\item[] \choice{D} Square of a sum
\end{itemize}
\end{minipage}

Work:

\vspace{1.4cm}

\noindent\rule{\textwidth}{0.4pt}

\vspace{6pt}
\textbf{4.} Use Pascal's Triangle to expand the expression $(x+2)^8$.

\begin{itemize}[leftmargin=1.2em, itemsep=6pt]
\item[] \choice{A} $x^8+256$
\item[] \choice{B} $256+x+16x^2+112x^3+448x^4+1120x^5+1792x^6+1792x^7+1024x^8$
\item[] \choice{C} $x^8+16x^7+112x^6+448x^5+1120x^4+1792x^3+1792x^2+1024x+256$
\item[] \choice{D} $x^7+14x^6+84x^5+280x^4+560x^3+672x^2+448x+128$
\end{itemize}

Show your work. Row 8 of Pascal's Triangle: \rule{8cm}{0.4pt}

Work:

\vspace{2.6cm}

\noindent\rule{\textwidth}{0.4pt}

\vspace{6pt}
\textbf{5.} Expand $(a+4)^5$:

\vspace{8pt}
\begin{center}
$a^5+\rule{1.3cm}{0.4pt}\,a^4+\rule{1.3cm}{0.4pt}\,a^3+\rule{1.3cm}{0.4pt}\,a^2+\rule{1.3cm}{0.4pt}\,a+\rule{1.3cm}{0.4pt}$
\end{center}

Work:

\vfill
\noindent\textbf{Scoring:} Question 4 and Question 5 are 10 points each. Question 2 is a bonus of up to 5 points. Total out of 20.

\newpage
\noindent{\LARGE\textbf{3-3 Lesson Quiz, short form: answer key}}

\vspace{4pt}\noindent\rule{\textwidth}{0.6pt}

\begin{enumerate}
\item[\textbf{2.}] \textbf{Part A D, $25-16x^6$; Part B A, Difference of squares.} $(5-4x^3)(5+4x^3)=5^2-(4x^3)^2=25-16x^6$. Choice C squares the variable but not the coefficient 4 (the TE's Try It 2 common error); A also multiplies the exponents $3\cdot3$; B is $(5-4x^3)^2$, a square, not a product of a sum and a difference. Plain distribution (Lesson 3-2) gets the same product. Savvas quiz item 2, DOK 2, Skills Review and Practice G100. Taught by TE Example 2 (cut in the fast packet; the packet's reference box names the identity). \textbf{Bonus 5:} 3 for Part A, 2 for Part B.
\item[\textbf{4.}] \textbf{C.} Row 8 is $1,8,28,56,70,56,28,8,1$; with $a=x$, $b=2$ the terms are $x^8$, $8(2)x^7=16x^7$, $28(4)x^6=112x^6$, $56(8)x^5=448x^5$, $70(16)x^4=1120x^4$, $56(32)x^3=1792x^3$, $28(64)x^2=1792x^2$, $8(128)x=1024x$, $2^8=256$. Choice A keeps only the first and last terms; B attaches the coefficients to the wrong powers of $x$; D is $(x+2)^7$ (Row 7). Savvas quiz item 4, DOK 2, G100. Taught by Example 5A and Try It 5a. \textbf{10 points:} 6 for C, 4 for a correct Row 8 or shown substitution.
\item[\textbf{5.}] \textbf{20; 160; 640; 1280; 1024.} Row 5 is $1,5,10,10,5,1$; $5(4)=20$, $10(4^2)=160$, $10(4^3)=640$, $5(4^4)=1280$, $4^5=1024$. Savvas quiz item 5, DOK 2, G100. Taught by Example 5B and Try It 5b. \textbf{10 points:} 2 per blank.
\end{enumerate}

{\emergencystretch=3em
\textbf{Items left off this form} (available for a retake or Form B): Savvas item 1, prove \mbox{$x^3-y^6$} $=$ \mbox{$(x-y^2)(x^2+xy^2+y^4)$} by filling in the four step reasons from a word bank (answer: Distributive; Multiply; $-1$; like terms; DOK 2, G100; TE Example 1, cut); Savvas item 3, factor $1-125n^3$, select all factors, and name the identity (answer: Part A A and E, $1-125n^3=(1-5n)(1+5n+25n^2)$; Part B B, Difference of cubes; DOK 2, G100; TE Example 3, cut). Item 3 is the same skill as Topic 3 assessment item 7 (factor $64+27a^3$) and IXL NJV (an optional, unstarred row in the packet's IXL table); if the PLC adds Example 3, move item 3 onto this form. The quiz has no graphs, so nothing was left off for a figure.\par}

\textbf{IXL replacement:} SmartScore 70 or above on A7M and NEU replaces this quiz with 20/20.

\end{document}
